% !PRINTED 437
\weberpage{437}
%<<CONT
and $\alpha$, $\beta$, $\gamma$ are the roots of the cubic equation
\[ \alpha^3 + \alpha^2 - \frac{1}{7} = 0. \]
By \S.\,72, (15), (12), (13) one may now set
\weq{15}{\Theta^2 = \tau^4\,\omega\,\tau^5\,\chi}
and from this one easily obtains, by (4) and (5),
\weq{16}{\Theta^2 = \begin{pmatrix} \beta, & \gamma\varepsilon^3, & \alpha\varepsilon^2 \\ \gamma\varepsilon^{-3}, & \alpha, & \beta\varepsilon^{-1} \\ \alpha\varepsilon^{-2}, & \beta\varepsilon, & \gamma \end{pmatrix}.}

If one compares this with the result that one obtains when $\Theta^2$ is formed directly from (10), one finds:
\weq{17}{\begin{aligned}
\alpha &= \alpha^2\varepsilon^{-1} + \beta^2\varepsilon + \gamma^2\varepsilon^{-2}, & \alpha &= \beta\gamma + \gamma\alpha\varepsilon^2 + \alpha\beta\varepsilon^{-1} \\
\beta &= \alpha^2\varepsilon^{-1} + \beta^2\varepsilon^3 + \gamma^2\varepsilon^{-3}, & \beta &= \beta\gamma\varepsilon^3 + \gamma\alpha + \alpha\beta\varepsilon \\
\gamma &= \alpha^2\varepsilon^2 + \beta^2\varepsilon^3 + \gamma^2\varepsilon^{-2}, & \gamma &= \beta\gamma\varepsilon^{-3} + \gamma\alpha\varepsilon^{-2} + \alpha\beta,
\end{aligned}}

and the comparison with the $\Theta^{-2}$ formed from (10) yields an entirely similar system of formulae, which arises from (17) when $\varepsilon$ is interchanged with $\varepsilon^{-1}$, namely:
\[ \begin{aligned}
\alpha &= \alpha^2\varepsilon + \beta^2\varepsilon^{-1} + \gamma^2\varepsilon^2, & \alpha &= \beta\gamma + \gamma\alpha\varepsilon^{-2} + \alpha\beta\varepsilon \\
\beta &= \alpha^2\varepsilon + \beta^2\varepsilon^{-3} + \gamma^2\varepsilon^3, & \beta &= \beta\gamma\varepsilon^{-3} + \gamma\alpha + \alpha\beta\varepsilon^{-1} \\
\gamma &= \alpha^2\varepsilon^{-2} + \beta^2\varepsilon^{-3} + \gamma^2\varepsilon^2, & \gamma &= \beta\gamma\varepsilon^3 + \gamma\alpha\varepsilon^2 + \alpha\beta.
\end{aligned} \]

But in order that no circular reasoning lie herein, it must be shown that the relations (17) are actually satisfied in consequence of the determinations (11), (12); for only then are the validity of the relation (15) and the equality of $\Theta^2$ with $\Theta^{-2}$ established. For this it suffices to derive two of these relations, say
\[ \begin{aligned}
\alpha &= \alpha^2\varepsilon^{-1} + \beta^2\varepsilon + \gamma^2\varepsilon^{-2} \\
\alpha &= \beta\gamma + \gamma\alpha\varepsilon^2 + \alpha\beta\varepsilon^{-1},
\end{aligned} \]
for once these have been established, one may interchange $\varepsilon$ with $\varepsilon^{-1}$ in them, whereby $\alpha$, $\beta$, $\gamma$ are not changed, and then $\varepsilon$ with $\varepsilon^2$ and with $\varepsilon^4$, whereby $\alpha$, $\beta$, $\gamma$ pass into $\gamma$, $\alpha$, $\beta$ and into $\beta$, $\gamma$, $\alpha$; and in this way one obtains the whole system of formulae (17).

These two relations, however, result from a simple calculation by (11), (12) in the form:
\[ \begin{aligned}
&-(\varepsilon + \varepsilon^2 + \varepsilon^4 - \varepsilon^{-1} - \varepsilon^{-2} - \varepsilon^{-4})(\varepsilon^4 - \varepsilon^{-4}) \\
&= \varepsilon^{-1}(\varepsilon^4 - \varepsilon^{-4})^2 + \varepsilon(\varepsilon^2 - \varepsilon^{-2})^2 + \varepsilon^{-2}(\varepsilon - \varepsilon^{-1})^2 \\
&= (\varepsilon - \varepsilon^{-1})(\varepsilon^2 - \varepsilon^{-2}) + \varepsilon^2(\varepsilon - \varepsilon^{-1})(\varepsilon^4 - \varepsilon^{-4}) \\
&\phantom{=} + \varepsilon^{-1}(\varepsilon^2 - \varepsilon^{-2})(\varepsilon^4 - \varepsilon^{-4}).
\end{aligned} \]
